% !TeX program = latexmk
% Speaker cue cards: one card environment per card, in speaking order. Each card
% gets the TU/e logo on one side and its text on the other:
%   \begin{card}[cue under the logo, optional]{heading} text \end{card}
% Print: double-sided, flip on long edge, A4, actual size (100 %); cut each
% sheet in half at the marks. Printer only flips on the short edge: use
% \usepackage[flip=short]{tuecards}.
% Too much text on a card is an error: shorten it or split it over two cards.
% Build: make (or latexmk -pdf cuecards.tex). Style: tuecards.sty.
\documentclass[a4paper]{article}
\usepackage{tuecards}

\reporttitle{Title of the talk}
\reportauthors{Name Surname}              % PDF author
% \cardlogo{figures/partner-logo.pdf}     % other logo (width: [86.44mm])
% \cardlogo{}                             % no logo
% \cardfontsize{12}{15}                   % smaller text (default 14/18 pt)
% \accentcolor{HTML}{003399}              % secondary color for headings etc.; default Scarlet

\begin{document}

\begin{card}{Opening (slide 1)}
  Good morning. I am Name Surname, and in the next ten minutes I will show you how our launcher reaches \SI{10}{\metre}.

  Start with the problem: the requirement, and why it is hard.
\end{card}

\begin{card}[spring -- latch -- 250 N/m]{The design (slide 2)}
  Point at the spring.
  The spring stores the energy; the latch releases it in one go.

  Key number: \SI{250}{\newton\per\metre}, compressed \SI{40}{\milli\metre}.
\end{card}

\begin{card}{Results (slides 3--4)}
  Three runs, all within \SI{0.5}{\metre} of the prediction.

  Pause. Let the graph speak.
\end{card}

\end{document}
